% AI-assisted working draft for team review, not an approved final section.
% Requires amsmath. Cites nothing.
\section{Conclusion}
\label{sec:conclusion}

We asked how geometry, supports and loads shape the deformation of a
two-dimensional elastic object. A hole makes the plate softer, and holes and
inner corners concentrate stress; moving the load on the bracket moves the
stress peak. Holding fewer displacement components barely changes the
displacement but raises $\kappa_2(A)$. Slender bending is the hardest case:
linear triangles are too stiff in bending, and even the finest beam mesh
deflects $0.7\%$ less than beam theory with shear predicts.

Since $A$ is sparse, symmetric and positive definite, sparse direct
factorization and CG both apply, and the CG bound predicts iterations
growing like $\sqrt{\kappa_2(A)}$. We measured $\kappa_2(A)$ growing like $N$ and CG
iterations like $N^{1/2}$ for every object; shape changes the constant, and
the beam is by far the worst conditioned. CG is the best iterative method.
Diagonal preconditioning saves few iterations and no time, Jacobi diverges
and Gauss--Seidel is impractically slow. On the beam MINRES stalls above the
residual target although its solution is accurate. At these sizes sparse LU
is the fastest method on every system.

Under refinement the displacements settle, but peak stresses do not: at the
inner corner the peak keeps growing (119 to 178~MPa), as the infinite exact
corner stress predicts. The algebraic error of the accepted CG solutions,
at most $10^{-9}$, is far below the mesh changes, so the mesh limits the
accuracy. The solutions satisfy the supports exactly and balance force and
moment, a load perturbation stays within the condition-number bound, and
removing the vertical support makes $A$ numerically singular, as expected. The main limitations, detailed in
Section~\ref{sec:results}, are that only diagonal preconditioning was
tested, constant element stresses make peaks mesh dependent, and the linear
model ignores yielding. Natural next steps are incomplete Cholesky or
multigrid preconditioners, larger or three-dimensional problems in which
iterative methods may overtake sparse LU, and better stress recovery or
higher-order elements.
